\section{Estructura categorial del TPK: estados enriquecidos, morfismos de
transporte y clasificación operatoria}
\label{sec:arbol-operatorio-tpk-rev11}
\index{TPK!clasificación operatoria}
\index{selector de fase!suboperador del TPK}

El TPK es el contenedor categórico que transporta sobre APP los estados
orientados de TRIT. Su definición precede al inventario de realizaciones:
primero se fijan objetos, morfismos, composición y clases operatorias; sólo
después se introducen los operadores que realizan cada clase.

\HMTClaim{HMT-R-TPK-ENRICHED-STATE}
\begin{definition}[TPK como categoría sobre la base observable con familia de fibras]
\label{def:estado-enriquecido-tpk-rev11}
Sea \(\mathsf P_{\rm obs}\) la categoría de rutas observables sobre el grafo
de Cayley de APP. Para cada configuración \(q\) se fija la fibra
\[
 \mathcal F_q=
 \mathsf O_q\times\mathsf H_q\times\mathsf C_q\times
 \mathsf S_q\times\mathsf B_q\times\mathsf\Lambda_q,
\]
cuyos factores conservan, respectivamente, orientación y hoja; residuo y
cociente; acarreo; firma; cilindro y frontera; y registro de transporte. Un
objeto del TPK es un par \(x=(q,f)\), con \(f\in\mathcal F_q\). Un morfismo
\(\Gamma:x\to y\) es una ruta \(\gamma:q\to q'\) de
\(\mathsf P_{\rm obs}\) junto con el transporte de fibra que preserva las
relaciones de extensión declaradas entre \(f\) y \(f'\). La identidad es la
ruta vacía con transporte identidad y la composición es la concatenación de
rutas con composición de los transportes de fibra. La proyección
\[
 p:\mathcal X_{\rm TPK}^{\rm enr}\longrightarrow\mathsf P_{\rm obs}
\]
envía cada estado a su configuración observable y cada morfismo a la ruta que
transporta. Esta definición fija una categoría sobre la base y una familia de
fibras; no presupone la existencia de levantamientos cartesianos para toda
flecha de \(\mathsf P_{\rm obs}\).
\end{definition}

\begin{definition}[Clasificación ontológico--operatoria del TPK]
\label{def:categorias-operatorias-tpk-rev11}
La estructura interna del TPK se organiza en ocho clases diferenciadas por tipo de
dominio, codominio y función:
\begin{enumerate}
 \item \textbf{Soporte y estado.} La base es \(\mathsf P_{\rm obs}\) y el
 portador es \(\mathcal X_{\rm TPK}^{\rm enr}\to\mathsf P_{\rm obs}\). Esta
 clase contiene objetos; las siete siguientes contienen mapas entre objetos
 o lectores de esos objetos.
 \item \textbf{Selección interna.} La función de fase
 \[
 M_{\rm ph}:\{1,\ldots,9\}\longrightarrow\{+1,-1,0\}
 \]
 produce el valor escalar que el operador
 \(\operatorname{Sel}_t\) recibe para activar la evaluación aditiva, la
 multiplicativa o el paso de umbral. Es \(\operatorname{Sel}_t\), y no
 \(M_{\rm ph}\) aisladamente, quien actúa sobre la coordenada de hoja; ninguno
 de los dos modifica por sí mismo la posición en APP.
 \item \textbf{Transporte.} Para cada dirección cardinal \(d\),
 \(T_d:X_9\to X_9\), \(p\mapsto p+\delta(d)\), transporta la posición. La
 evolución \(F_{\rm obs}:\mathcal Q_{\rm obs}\to\mathcal Q_{\rm obs}\)
 compone selección, traslación, orientación y avance de fase.
 \item \textbf{Lectores y cocientes.} \(\rho_9\) publica la clase nonádica;
 \((r_3,c_3)\) conserva residuo trítico y cociente; y
 \((q_{1000},r_{1000})\) separa bloque decimal y acarreo. Son mapas de
 lectura, no transiciones del estado.
 \item \textbf{Memoria y frontera.} Los transportes de
 \(\mathsf H_q,\mathsf C_q,\mathsf S_q,\mathsf B_q\) y
 \(\mathsf\Lambda_q\) actualizan cociente, acarreo, firma, cilindro, frontera
 y registro. Su proyección visible puede coincidir aunque los estados
 enriquecidos difieran.
 \item \textbf{Periodicidad y retorno.} La extensión
 \[
 0\longrightarrow C_{12}\longrightarrow C_{108}
 \longrightarrow C_9\longrightarrow0
 \]
 tipa sector, calendario y fase visible. El retorno nonádico restaura la
 fase; la periodicidad de ciento ocho iteraciones restaura el calendario;
 ninguno borra
 por definición la memoria de frontera. El lazo
 \(\mathcal K_{27}=F_{\rm obs}^{27}\) es un morfismo compuesto, no otro
 reloj.
 \item \textbf{Prolongación.} Las relaciones
 \(\operatorname{Ext}_r\subseteq\mathcal F_q\times\mathcal F_{q'}\)
 seleccionan continuaciones compatibles. Los cilindros y su sistema inverso
 publican supervivencia finita y, bajo las hipótesis de compatibilidad
 declaradas, una sección global.
 \item \textbf{Salidas.} Los funtores de evaluación contraen el mismo estado
 enriquecido hacia la salida arquimediana, la salida de incidencia
 excepcional o los clasificadores del atlas. Estas salidas son codominios de
 lectura; no constituyen nuevos contenedores junto a TPK.
\end{enumerate}
\end{definition}

Cada operador que aparece a continuación queda determinado por la séxtupla
\((D,C,f,I,L,M)\): dominio \(D\), codominio \(C\), regla \(f\), invariantes
preservados \(I\), información publicada o perdida \(L\) y memoria conservada
\(M\). La igualdad de cardinales, periodos o nombres históricos no identifica
dos operadores con séxtuplas distintas. Una vez fijada esta
clasificación, el siguiente diagrama resume cinco agrupaciones funcionales
de realizaciones que refinan las ocho clases operatorias del TPK:
\[
\begin{tikzpicture}[
  node distance=7mm and 9mm,
  every node/.style={font=\small,align=center},
  root/.style={draw=hmtblue,very thick,rounded corners,fill=hmtlightblue,
    inner sep=5pt},
  op/.style={draw=hmtblue,rounded corners,fill=white,inner sep=4pt},
  sub/.style={draw=hmtgray,rounded corners,fill=hmtpaper,inner sep=3pt},
  edge/.style={-{Latex[length=2mm]},semithick,hmtblue}
]
\node[root] (tpk) {TPK\\$\mathcal X_{\rm TPK}^{\rm enr}$};
\node[op,below left=of tpk] (sel) {selección\\de lectura};
\node[op,below=of tpk] (trans) {transporte\\y retorno};
\node[op,below right=of tpk] (lift) {elevación\\y memoria};
\node[sub,below=of sel] (mov) {$\operatorname{Sel}_t[\,M_{\rm ph}(g(t))\,]$};
\node[sub,below=of trans] (card) {$T_d,\ \mathcal K_{27},\ F_{\rm obs}$};
\node[sub,below=of lift] (fib) {$\mathsf H,\mathsf C,\mathsf S,
  \mathsf B,\mathsf\Lambda$};
\node[op,below=17mm of trans] (clock) {calendario y registro\\
  $C_{12}\hookrightarrow C_{108}\twoheadrightarrow C_9$, $\Theta_{108}$};
\node[op,below=of clock] (ext) {prolongación, cilindros, frontera\\
  y filtración de supervivencia};
\draw[edge] (tpk) -- (sel);
\draw[edge] (tpk) -- (trans);
\draw[edge] (tpk) -- (lift);
\draw[edge] (sel) -- (mov);
\draw[edge] (trans) -- (card);
\draw[edge] (lift) -- (fib);
\draw[edge] (mov) |- (clock);
\draw[edge] (card) -- (clock);
\draw[edge] (fib) |- (clock);
\draw[edge] (clock) -- (ext);
\end{tikzpicture}
\]
El diagrama fija una contención, no una sucesión de teorías. En particular,
el selector interno, las traslaciones cardinales, la observación
estroboscópica y los lectores modulares conservan los tipos fijados en la
\cref{def:categorias-operatorias-tpk-rev11}. El TPK coordina sus composiciones
y conserva su procedencia en un mismo estado enriquecido.

\HMTClaim{HMT-R-TPK-OPERATOR-SEPARATION}
\begin{proposition}[Separación de selección, movimiento y observación]
\label{prop:tpk-separacion-operadores-rev11}
La selección \(\operatorname{Sel}_t[\,M_{\rm ph}(g(t))\,]\), el transporte
cardinal \(\operatorname{Lift}(T_d)\) y el muestreo
\(\operatorname{Obs}_3\) son operaciones no intercambiables: la primera elige
la evaluación activa; la segunda modifica la posición y el enrollamiento; la
tercera publica una subsucesión temporal. Su composición dentro del TPK
determina una ruta observable, mientras la fibra \(\mathcal F_q\) conserva
los datos necesarios para distinguir prolongaciones con igual extremo.
\end{proposition}

\begin{proof}
Los tipos son diferentes: \(M_{\rm ph}\) toma valores en el conjunto ternario
de activación, \(\operatorname{Sel}_t\) actualiza la coordenada de hoja,
\(\operatorname{Lift}(T_d)\) transporta el estado sobre el toro y
\(\operatorname{Obs}_3\) actúa sobre sucesiones. Una identificación entre dos
de estos objetos violaría el tipado de dominio o codominio. La extensión por fibras
incorpora, además, residuo--cociente, acarreo, frontera y registro, de modo
que la proyección al extremo del camino es un mapa de olvido explícito.
\end{proof}

\begin{statusbox}[title={Nomenclatura canónica}]
La denominación pública del contenedor es \emph{TPK}. Cuando sea necesario
nombrar el dato transportado se empleará \emph{estado enriquecido del TPK} o
\(\mathcal X_{\rm TPK}^{\rm enr}\). La función \(M_{\rm ph}\) designa
exclusivamente el valor interno de fase que recibe
\(\operatorname{Sel}_t\).
\end{statusbox}
